\documentclass{resume}
\usepackage{zh_CN-Adobefonts_external} 
\usepackage{linespacing_fix}
\usepackage{cite}
\usepackage{comment}

\usepackage[English]{languageSelection}

\usepackage[colorlinks=true, linkcolor=blue, urlcolor=blue]{hyperref}

\usepackage{noteplus} 

\begin{document}
\pagenumbering{gobble}

% Introduction
\EN{
  \name{赵禹惟}
}

\EN{
  \begin{minipage}[t]{0.92\textwidth}
    \info{本科生}{中国海洋大学}{Heriot-Watt Univerisity}{}
  \end{minipage}

  \vspace{0.5em}

  \info{电话: (+86) 13578774880}{}{}{}
  \info{邮件: stanley\_zhao0113@outlook.com}{}{}{}
  \info{个人主页: https://github.com/RamessesN}{}{}{}
}

\yourphoto{0.15}

% 个人简介
\section{个人简介}\vspace{0.5em}
我是赵禹惟，来自吉林省长春市，本科阶段就读于 \textbf{中国海洋大学} 与 \textbf{英国赫瑞瓦特大学} 联合办学项目，主修计算机科学和机器人双学位。

在校期间，我积极参与各类项目与竞赛，积累了丰富的实践经验，尤其是在嵌入式系统开发、计算机视觉和机器人控制等领域。
我具备扎实的编程能力与系统实现经验，熟悉 STM32、Jetson Orin、RealityKit、TensorRT 等软硬件平台。在多项科研与工程项目中，我担任核心开发者，展现出良好的团队协作能力与问题解决能力。

未来，我希望进一步深入研究人工智能与机器人技术的融合方向 - \textbf{“具身智能”}，致力于推动该领域的技术创新与实际应用发展。

\vspace{0.8em}

% 教育背景
\section{教育背景}\vspace{0.5em}
\begin{itemize}

  \item 
  \begin{tabular*}{0.95\textwidth}{@{\extracolsep{\fill}}lr}
    \textbf{中国海洋大学} & \textit{2023.08 -- 2026.06} \\
    \textit{计算机科学学士学位} & \\
  \end{tabular*}

  \vspace{0.8em}

  \item 
  \begin{tabular*}{0.95\textwidth}{@{\extracolsep{\fill}}lr}
    \textbf{Heriot-Watt University} & \textit{2026.09 -- 2027.06} \\
    \textit{机器人工学学士学位} & \\
  \end{tabular*}

  \vspace{1em}

  \textbf{[核心课程]} 计算机系统基础 (98.5 / 100), 人工智能 (96.5 / 100), 机器人集成项目 (96.5 / 100), Linux \& C (96 / 100), 线性代数 (95 / 100)

\end{itemize}

\vspace{0.8em}

% 科研和项目经历
\section{科研和项目经历}\vspace{0.5em}
\begin{itemize}
  \item \textbf{科研经历}
  \begin{itemize}
    \item \textbf{OverLoCK 特征提取与多源遥感图像分类研究}·2025-2026·推进中·指导老师: 高峰\,[OUC]
  \end{itemize}

  \item \textbf{Robotics Group Project I\;课程助教}·2025·指导老师: 张述\,[OUC]
  \begin{itemize}
    \item 设计了实验流程和测试环境，协助批改了 90+ 学生的作业和实验报告。
    \item 在实验室课程中讲解了 Python 语法和 Robomaster SDK 的使用方法，帮助学生解决编程和硬件方面的问题。
    \item 将项目平均完成率提高了 20\%。
  \end{itemize}
  
  \item \textbf{项目经历}
  \begin{itemize}
    \item \textbf{基于 Robomaster-SDK 实现对 Python 3.8 及以上版本的兼容性支持}·2025· \href{https://github.com/RamessesN/Robomaster-SDK-Ultra}{Github @Robomaster-SDK-Ultra}
    \item \textbf{基于 Swift 重写路径追踪算法并以取代传统的光栅化技术用于黑洞模拟实验}·2025· \href{https://github.com/RamessesN/Black_Hole_Simulator}{GitHub @BlackHole\_Simulator}
    \item \textbf{海上光伏监测与运行系统}·2025·系统构建与信号传输；网络智能控制平台建设·推进中
    \item \textbf{多模态融合：全天候海上智能感知}·2025·利用单目深度算法生成雷达点云图像，以补充数据集· \href{https://github.com/RamessesN/VesselContest_C3}{GitHub @VesselContest\_C3}
    \item \textbf{Swift 原生 iOS 增强现实应用程序开发}·2025·基于 RealityKit 和 ARKit 实现任意空间锚点设置及导航技术· \href{https://github.com/RamessesN/Macintustin}{GitHub @Macintustin}
    \item \textbf{将 Nim 语言移植至龙芯架构 LoongArch64 上}·2024·在基于 LLVM 后端的龙芯架构上实现尼姆的环境适应功能· \href{https://github.com/RamessesN/Nim2LoongArch64}{GitHub @Nim2LoongArch64}
  \end{itemize}
\end{itemize}

\vspace{0.8em}

% 荣誉和奖励
\section{荣誉和获奖}\vspace{0.5em}
\begin{itemize}
  \item 2025·\textbf{全国海洋航行器设计与制作大赛 \textit{C3}}·全国二等奖·负责利用单目深度算法生成雷达点云图像，以补充数据集
  \item 2025·\textbf{ICM 美国大学生数学建模大赛}·Honorable Mention·负责数学模型的构建以及数据挖掘工作
  \item 2025·\textbf{中国高校计算机大赛 - 移动应用创新赛}·华东赛区二等奖·负责 Swift 语言编写的 iOS 应用程序的整个开发流程
  \item 2024·\textbf{全国大学生计算机系统能力大赛}·华东赛区优胜奖·负责 Nim 在 LoongArch64 架构上的环境适配调整工作
  \item 2025·\textbf{全国海洋航行器设计与制作大赛 \textit{F1}}·齐鲁赛区三等奖·负责 STM 硬件的构建以及算法的实现
  \item 2024·\textbf{Apple Swift Certification Contest}·Certified User
  \item 2024-2025·\textbf{中国海洋大学校级奖学金}·综合三等奖学金 (9 / 48)
  \item 2023-2024·\textbf{中国海洋大学校级奖学金}·综合三等奖学金 (12 / 47)
\end{itemize}

\vspace{0.8em}

% 技术栈与工具
\section{技术栈与工具}\vspace{0.5em}
\begin{itemize}
  \item \textbf{编程语言: }C/C++, Java, Python, Swift, HTML/CSS, JavaScript, Nim
  \item \textbf{开发框架\&库: }RealityKit, ARKit, OpenCV, YOLOv8, STM32 HAL, TensorRT, PyTorch
  \item \textbf{嵌入式系统: }STM32, Arduino, ESP32, Jetson Orin, PWM control, TIM interrupt, sensor integration
  \item \textbf{开发平台: }Linux(Ubuntu), macOS, CUDA GPU acceleration, ROS2(basic)
  \item \textbf{开发工具: }Git, Docker, Xcode, VSCode, CMake, LaTeX, Markdown
  \item \textbf{外语: }英语(CET-6, 520), 德语(基础交流)
\end{itemize}

\end{document}